\textbf{Phần 1. Trắc nghiệm lựa chọn}
\setcounterpageref{ex}{0}

\begin{ex} %[2D8V41]
Có $12$ quả cam chia đều vào $2$ đĩa. Mỗi đĩa có bao nhiêu quả cam?
\choice
{$5$ quả cam}
{$7$ quả cam}
{\True $6$ quả cam}
{$4$ quả cam}
\loigiai{
$12 : 2 = 6$ (quả cam).
\textbf{Đáp án của câu hỏi là C}
}
\end{ex}

\begin{ex} %[2D8V41]
Từ phép nhân $2 \times 6 = 12$, ta viết được phép chia nào dưới đây?
\choice
{$12 : 3 = 4$}
{\True $12 : 2 = 6$}
{$6 : 2 = 3$}
{$12 : 6 = 3$}
\loigiai{
Từ phép nhân $2 \times 6 = 12$ ta có hai phép chia: $12 : 2 = 6$ và $12 : 6 = 2$.
\textbf{Đáp án của câu hỏi là B}
}
\end{ex}

\begin{ex} %[2D8V41]
Có $15$ bông hoa cắm đều vào các lọ, mỗi lọ cắm $5$ bông hoa. Hỏi cắm được bao nhiêu lọ hoa?
\choice
{$4$ lọ hoa}
{$2$ lọ hoa}
{$5$ lọ hoa}
{\True $3$ lọ hoa}
\loigiai{
$15 : 5 = 3$ (lọ hoa).
\textbf{Đáp án của câu hỏi là D}
}
\end{ex}

\begin{ex} %[2D8V41]
Kết quả của phép chia $20 : 2$ là
\choice
{\True $10$}
{$8$}
{$9$}
{$12$}
\loigiai{
$20 : 2 = 10$.
\textbf{Đáp án của câu hỏi là A}
}
\end{ex}

\begin{ex} %[2D8V41]
Một sợi dây dài $18\text{ dm}$ được cắt thành các đoạn bằng nhau, mỗi đoạn dài $2\text{ dm}$. Cắt được bao nhiêu đoạn dây?
\choice
{$8$ đoạn dây}
{$10$ đoạn dây}
{\True $9$ đoạn dây}
{$7$ đoạn dây}
\loigiai{
$18 : 2 = 9$ (đoạn dây).
\textbf{Đáp án của câu hỏi là C}
}
\end{ex}

\begin{ex} %[2D8V41]
Tính giá trị biểu thức: $14 : 2 + 18$.
\choice
{$24$}
{\True $25$}
{$26$}
{$23$}
\loigiai{
$14 : 2 = 7$; $7 + 18 = 25$.
\textbf{Đáp án của câu hỏi là B}
}
\end{ex}

\begin{ex} %[2D8V41]
Có $30$ bạn học sinh chia đều thành các nhóm, mỗi nhóm có $5$ bạn. Hỏi chia được bao nhiêu nhóm?
\choice
{$5$ nhóm}
{$7$ nhóm}
{$8$ nhóm}
{\True $6$ nhóm}
\loigiai{
$30 : 5 = 6$ (nhóm).
\textbf{Đáp án của câu hỏi là D}
}
\end{ex}

\begin{ex} %[2D8V41]
Điền dấu thích hợp: $16 : 2 \dots 25 : 5$.
\choice
{\True $>$}
{$<$}
{$=$}
{Không so sánh được}
\loigiai{
$16 : 2 = 8$; $25 : 5 = 5$. Vì $8 > 5$ nên điền dấu $>$. 
\textbf{Đáp án của câu hỏi là A}
}
\end{ex}

\begin{ex} %[2D8V41]
Có $10$ chiếc bút chì chia đều cho Mai và Nam. Mỗi bạn nhận được bao nhiêu chiếc bút chì?
\choice
{$4$ chiếc}
{$6$ chiếc}
{\True $5$ chiếc}
{$3$ chiếc}
\loigiai{
$10 : 2 = 5$ (chiếc).
\textbf{Đáp án của câu hỏi là C}
}
\end{ex}

\begin{ex} %[2D8V41]
Một hộp bánh có $20$ chiếc bánh. Nếu chia đều cho $5$ bạn thì mỗi bạn được mấy chiếc bánh?
\choice
{$3$ chiếc bánh}
{\True $4$ chiếc bánh}
{$5$ chiếc bánh}
{$6$ chiếc bánh}
\loigiai{
$20 : 5 = 4$ (chiếc bánh).
\textbf{Đáp án của câu hỏi là B}
}
\end{ex}

\begin{ex} %[2D8V41]
Tìm $x$ biết: $x \times 2 = 18$.
\choice
{$8$}
{$7$}
{$10$}
{\True $9$}
\loigiai{
$x = 18 : 2 = 9$.
\textbf{Đáp án của câu hỏi là D}
}
\end{ex}

\begin{ex} %[2D8V41]
Một thanh gỗ dài $16\text{ m}$ cưa thành các khúc gỗ dài $2\text{ m}$. Hỏi phải thực hiện bao nhiêu nhát cưa?
\choice
{\True $7$ nhát cưa}
{$8$ nhát cưa}
{$6$ nhát cưa}
{$9$ nhát cưa}
\loigiai{
Số khúc gỗ: $16 : 2 = 8$ (khúc). Số nhát cưa: $8 - 1 = 7$ (nhát cưa).
\textbf{Đáp án của câu hỏi là A}
}
\end{ex}

\textbf{Phần 2. Trắc nghiệm đúng sai}
\setcounterpageref{ex}{0}

\begin{ex} %[2D8V41]
Xét tính đúng sai của các nhận định về phép chia:
\choiceTF
{\True Từ một phép nhân hai thừa số khác nhau có thể lập được hai phép chia tương ứng.}
{Số 0 có thể làm số chia trong phép chia.}
{\True Một số tự nhiên khác 0 chia cho chính nó luôn bằng 1.}
{Phép tính 18 : 2 có kết quả là 8.}
\loigiai{
\begin{itemize}
    \item Ý a) Đúng: Đúng: $a \times b = c \Rightarrow c : a = b$ và $c : b = a$.
    \item Ý b) Sai: Không thể chia cho số 0.
    \item Ý c) Đúng: $a : a = 1$. Đúng.
    \item Ý d) Sai: $18 : 2 = 9$ (không phải 8).
\end{itemize}
\begin{itemize}
    \item a) \textbf{Đúng}.
    \item b) \textbf{Sai}.
    \item c) \textbf{Đúng}.
    \item d) \textbf{Sai}.
\end{itemize}
}
\end{ex}

\begin{ex} %[2D8V41]
Có 20 cái kẹo chia đều cho 5 bạn:
\choiceTF
{\True Mỗi bạn nhận được 4 cái kẹo.}
{\True Nếu có thêm 5 cái kẹo nữa thì mỗi bạn nhận được 5 cái kẹo.}
{Nếu chia đều cho 2 bạn thì mỗi bạn được 8 cái kẹo.}
{\True Số kẹo của mỗi bạn bằng một phần năm tổng số kẹo.}
\loigiai{
\begin{itemize}
    \item Ý a) Đúng: $20 : 5 = 4$ cái kẹo.
    \item Ý b) Đúng: $25 : 5 = 5$ cái kẹo.
    \item Ý c) Sai: $20 : 2 = 10$ cái kẹo (không phải 8).
    \item Ý d) Đúng: Đúng.
\end{itemize}
\begin{itemize}
    \item a) \textbf{Đúng}.
    \item b) \textbf{Đúng}.
    \item c) \textbf{Sai}.
    \item d) \textbf{Đúng}.
\end{itemize}
}
\end{ex}

\textbf{Phần 3. Trả lời ngắn}
\setcounterpageref{ex}{0}

\begin{ex} %[2D8V41]
Tính kết quả: 18 : 2.
\shortans{9}
\loigiai{
\textbf{Đáp số:} $9$\\
$18 : 2 = 9$.
}
\end{ex}

\begin{ex} %[2D8V41]
Có 25 lít dầu chia đều vào 5 can. Hỏi mỗi can chứa bao nhiêu lít dầu?
\shortans{5}
\loigiai{
\textbf{Đáp số:} $5$\\
$25 : 5 = 5\text{ L}$.
}
\end{ex}

\begin{ex} %[2D8V41]
Tính giá trị: 16 : 2 + 35.
\shortans{43}
\loigiai{
\textbf{Đáp số:} $43$\\
$16 : 2 = 8$; $8 + 35 = 43$.
}
\end{ex}

\begin{ex} %[2D8V41]
Một cuộn dây dài 14 m được cắt thành các đoạn dài 2 m. Hỏi cắt được bao nhiêu đoạn dây?
\shortans{7}
\loigiai{
\textbf{Đáp số:} $7$\\
$14 : 2 = 7$ đoạn.
}
\end{ex}

\begin{ex} %[2D8V41]
Tìm x biết: x x 5 = 35.
\shortans{7}
\loigiai{
\textbf{Đáp số:} $7$\\
$x = 35 : 5 = 7$.
}
\end{ex}

\begin{ex} %[2D8V41]
Có 20 kg gạo chia đều vào 2 bao. Hỏi mỗi bao có bao nhiêu ki-lô-gam gạo?
\shortans{10}
\loigiai{
\textbf{Đáp số:} $10$\\
$20 : 2 = 10\text{ kg}$.
}
\end{ex}
